\clearpage\thispagestyle{empty}
\pdfbookmark[0]{Quellen zum Pallas-Bericht vom Oktober 1802}{pallas180210-sources}
\begin{center}\textbf{Quellen zum Pallas-Bericht vom Oktober 1802}\end{center}
\noindent Diese KI-Übertragung in heutiges Deutsch folgt dem Auszug aus Franz Xaver von Zachs Bericht über Pallas in \href{https://gdz.sub.uni-goettingen.de/id/PPN236007467}{Gauß’ \textit{Werken}, Band~VI} (Göttingen, 1874), S.~226--227. Von Zach beschreibt Gauß’ vierte Korrektur der Bahnelemente und die daraus berechnete Ephemeride. Der Auszug enthält sechs Werte für die Lichtstärke und den Vergleich zweier vorausberechneter Positionen.
\par\smallskip\noindent Der \href{https://api.digitale-sammlungen.de/iiif/image/v2/bsb10538598_00400/full/1450,/0/default.jpg}{Erstdruck vom Oktober 1802} steht als Beitrag XXXIX in der \textit{Monatlichen Correspondenz}, Band~VI, S.~390--396. Der Auszug in den \textit{Werken} reicht von \href{https://api.digitale-sammlungen.de/iiif/image/v2/bsb10538598_00403/full/1450,/0/default.jpg}{Gauß’ vierter Korrektur auf S.~393} bis zum \href{https://api.digitale-sammlungen.de/iiif/image/v2/bsb10538598_00406/full/1450,/0/default.jpg}{Vergleich und dem Absatz über die Opposition auf S.~396}.
\par\smallskip\noindent Die Zeitschrift druckt die \href{https://api.digitale-sammlungen.de/iiif/image/v2/bsb10538598_00404/full/1450,/0/default.jpg}{sieben Angaben zu den Elementen IV auf S.~394} und die \href{https://api.digitale-sammlungen.de/iiif/image/v2/bsb10538598_00405/full/1450,/0/default.jpg}{Ephemeride auf S.~395}. In den \textit{Werken} stehen ein Verweis auf den Septemberbrief an Bode und eine Beschreibung der Ephemeride in eckigen Klammern. Die Tabelle der Zeitschrift gibt für die tägliche Bewegung \(769{,}''726\) an und beendet die Angabe zum Aphel mit 41. Im angeführten Brief stehen \(769''7263\) und am Ende der Angabe zum Aphel \(42''\). Im Vergleich der \textit{Werke} kommen aus der Handschrift zwei Spalten für die Elemente V und am Schluss die Angabe der Opposition nach V hinzu.
\par\smallskip\noindent Übertragung und Satz wurden mit KI erstellt. Die Digitalisate der Zeitschrift stammen von der Bayerischen Staatsbibliothek.
\clearpage
